\section{Introduction}

The repeatability of a kinematic eye-movement marker is usually summarized by one intraclass correlation (ICC), yet that ICC is conditional on the pipeline that produced the events. Event-detection algorithms differ widely and systematically in the events they return~\citep{andersson2017,birawo2022,syscls2026}, researchers define fixations and saccades in different ways~\citep{hessels2018}, and video-based tracker signals contain post-saccadic oscillations whose size depends on the saccade~\citep{hooge2015}. Processing that discards segments of uncertain quality also changes derived estimates such as fixation durations~\citep{wass2013}. A single ICC therefore does not show how much of the reported repeatability depends on which events were admitted.

This dependence matters for binocular measures. The two eyes are not perfectly yoked during reading. Binocular coordination has been studied in reading and non-reading tasks~\citep{liversedge2006,kirkby2008}; reported effects include interocular differences in saccade amplitude and post-saccadic drift~\citep{vernet2009} and fixation disparity that varies within and between fixations~\citep{nuthmann2009}. A classic asymmetry concerns horizontal saccades: the abducting eye tends to move faster than the adducting eye~\citep{collewijn1988}. Interocular differences in saccade kinematics are of interest as individual-difference and clinical measures. Ratios of abducting to adducting peak velocity are part of a standardized saccade protocol~\citep{nijbijvank2018}, individual differences in the asymmetry of binocular saccades have been analysed with mixed-effects models~\citep{kloke2006}, and binocular saccade coordination has been studied in children with dyslexia~\citep{bucci2008,jainta2011} and, during reading, in children with strabismus~\citep{lions2013}. Eye tracking is increasingly integrated into virtual-reality head-mounted displays (HMDs)~\citep{adhanom2023} and has been used to assess saccades~\citep{imaoka2020}, and binocular HMD recordings now exist for hundreds of participants~\citep{lohr2023}, so such measures can be obtained from large samples. Binocular data add a construction choice. A detector applied to the mean of the two per-eye speed magnitudes (here, mean-speed detection) yields one event window per saccade, from which a peak velocity is read for each eye. A detector applied separately to each eye yields per-eye events that must be paired. In the data analysed here, some mean-speed events were not matched by a per-eye detection in both eyes, and their left/right peak-velocity ratios were much further from equality than those of supported events. Whether such events are admitted therefore changes which ratios enter an interocular marker.

We examined this dependence in the reading task of GazeBaseVR, a public, longitudinal, binocular HMD eye-tracking dataset~\citep{lohr2023}. The marker, $A$, contrasts the log left/right-eye peak-velocity ratio of leftward return sweeps with that of rightward progressive saccades within a recording. We computed $A$ under six event constructions, which differ in detector (mean-speed or per-eye), in admission (all mean-speed events, or only those overlapping a matched per-eye pair) and in an optional recording-centred abstention rule. We then compared their short-term ICCs on a common cohort with shared bootstrap draws, separately from the effect of recording eligibility. The approach follows the rule-based kinematic event definitions, internal rule-consistency audit and threshold-sensitivity reporting of our earlier virtual-reality work~\citep{vrst26}. The same dataset has been used for eye-movement biometrics, including binocular features~\citep{raju2024,yao2026}, and an earlier round-1 release from a long-term virtual-reality data collection by the same group was used for biometrics~\citep{lohr2020}. Here the question is how event construction moves the reliability estimate of one interocular marker.

On 376 participants eligible under all six constructions (eligibility floor 5), the ICC(A,1) between the first two sessions (R1 S1--S2, about 30 minutes apart) rose from .374 for unfiltered mean-speed events to .693 when those events were admitted only if a matched per-eye detection overlapped them (path-specific $\Delta$ICC .319, descriptive 95\% CI .003 to .501), with event windows, peaks and classes unchanged. Most of this difference came from one participant's extreme session difference under unfiltered events: omitting that participant from both constructions left a difference of .054 ($-$.007 to .129). Many of that session's return sweeps were physiologically implausible one-eye events that the unfiltered mean-speed detector accepted, so the contrast is best read as the leverage of a single session on a permissive baseline. A post hoc comparator that only removed events with a peak velocity above a conventional 1,000$^\circ$/s gave a similar ICC (.711 against .693; admitted minus ceiling $-$.018, $-$.090 to .039), and whether the unfiltered baseline was low at all depended on detector settings: across velocity thresholds, smoothing windows and pairing gaps its ICC ranged from .226 to .703, whereas admitted and ceiling ICCs stayed between .66 and .74 (each variant in its own eligible cohort). Admission is therefore one of several plausibility rules that limit the leverage of the extreme session, while removing largely different events, and showed no consistent advantage over a ceiling; in the units of $A$, the median absolute session difference was .019 unfiltered and .017 admitted. The analysis is exploratory and post hoc (Methods), and the admitted constructions were defined after the results were seen. We report path-specific contrasts on the common cohort alongside the effect of recording eligibility, which acted largely through the same few extreme sessions, locate where the constructions differ through influence and agreement analyses, and test what the abstention rule detects under synthetic one-eye faults, together with amplitude-matched left/right controls, repeatability across about two years and robustness to constant per-eye gain.
